\section{Discussion}
\label{sec:discussion}

\paragraph{Why the cross-level coupling design works.}
The three coupling channels of \cref{sec:tap} close a feedback loop that is absent when the two levels train independently.
Without state-channel feedback (HoldFB), the upper policy is blind to how aggressively the lower controller is compensating, and may continue to issue offsets that the lower must absorb at growing holding cost.
Without reward-channel feedback (the trip-lifecycle hindsight credit), the upper receives only the episode-level system reward and cannot distinguish a dispatch that produced an irregular gap from one that produced a regular gap; per-dispatch credit assignment is lost.
We deliberately keep the lower channel one-way: only the action channel transmits upper signal to the lower, and TPC (\cref{sec:tpc}) corrects for the resulting train-time covariate shift. We initially attempted reverse cross-advantage (lower reward = $r_{\mathrm{L}} + \beta R_{\mathrm{U}}$ for transitions in trip $i$, in the spirit of RoboDuet) but found that the high-variance per-trip $R_{\mathrm{U}}$ injected scale-mismatched noise into lower SAC training and degraded performance; the reward-channel-forward, action-channel-backward design we use is markedly more stable.

\paragraph{Connection to RoboDuet.}
TransitDuet draws inspiration from the RoboDuet framework of \citet{pan2024roboduet}, which introduced cross-advantage augmentation for bi-level robot locomotion control.
\Cref{tab:analogy} makes the structural analogy explicit.
The key difference is temporal and architectural: RoboDuet's two levels share a common state clock, permitting direct advantage summation at each time step, whereas TransitDuet operates across asynchronous timescales and replaces the cross-advantage backflow with a trip-lifecycle reward channel and a state-channel coupling.

\begin{table*}[!t]
\centering
\caption{Structural analogy between RoboDuet and TransitDuet.}
\label{tab:analogy}
\begin{tabular}{@{}lll@{}}
\toprule
\textbf{Aspect} & \textbf{RoboDuet} & \textbf{TransitDuet} \\
\midrule
Upper policy     & Locomotion style selector    & Timetable planner ($\piU$) \\
Lower policy     & Joint-level tracking controller & Holding controller ($\piL$) \\
Upper action     & Style parameters             & Per-dispatch offset $\delta_t \in [-120, 120]$\,s \\
Lower action     & Joint torques                & Holding duration (s) \\
Coupling         & Synchronous cross-advantage  & State channel (HoldFB) + reward channel (hindsight credit) + action channel ($\delta_t$) \\
Timescale        & Shared clock (same $\Delta t$) & Heterogeneous ($300$\,s vs.\ $1$\,s) \\
Constraint       & Joint-limit penalties        & Fleet ($\thetavec$-OGD) + headway (Lagrangian) \\
Activation schedule & Fixed $\beta$              & Lower-only warmup ($e \leq 30$) then full coupling \\
\bottomrule
\end{tabular}
\end{table*}

\paragraph{Limitations.}
Several limitations qualify the present results and motivate future work.

\emph{Single-route scope.}
TransitDuet is evaluated on a single bidirectional corridor.
Real transit networks involve multiple interacting routes with shared terminals, transfer passengers, and fleet rebalancing across lines.
Extending the framework to a multi-route setting would require the upper policy to allocate a shared fleet across routes, substantially expanding the action space and introducing inter-route coordination constraints.

\emph{Simplified demand model.}
Passenger arrivals follow stationary OD rates within each time period.
In practice, demand is non-stationary, exhibits day-of-week and seasonal variation, and responds endogenously to service quality (e.g., passengers abandoning a stop after excessive waits).
Incorporating demand elasticity and real-time demand forecasting would improve realism but introduce additional model complexity.

\emph{No real-world deployment.}
All experiments are conducted in a stochastic microsimulation.
While the simulator reproduces key phenomena---bunching, load-dependent dwell times, stochastic travel times---it abstracts away traffic signal interactions, driver behavior heterogeneity, and vehicle breakdowns.
Sim-to-real transfer would require domain randomization or a calibrated digital twin.

\emph{Computational cost.}
Each episode requires approximately 7\,s of environment simulation time plus training overhead.
At 300 episodes with 3 seeds, the total wall-clock time is approximately 14 hours per seed on a single 12\,GB GPU.
While acceptable for research, real-time retraining for schedule disruptions would require further acceleration, for instance through model distillation or warm-starting from a pre-trained policy.

\paragraph{Scalability considerations.}
Scaling to multi-route networks introduces two challenges.
First, the upper-level action space grows combinatorially with the number of routes, since headway targets must be set jointly to respect shared fleet constraints.
Factored action representations or attention-based architectures over route embeddings could mitigate this.
Second, the lower-level parameter sharing, which currently treats all buses on a single route identically, would need to incorporate route-specific context to generalize across heterogeneous lines.
Graph neural networks operating on the transit network topology offer a promising direction for encoding such structure.
